\documentclass[11pt]{article}
\usepackage[margin=1in]{geometry}
\usepackage{amsmath,amssymb,amsthm}
\usepackage[hidelinks]{hyperref}

\newtheorem{theorem}{Theorem}
\newcommand{\BPBp}{\mathrm{BPBp}}
\newcommand{\ot}{\widehat{\otimes}_{\varepsilon}}

\title{Historical partial packet: unital subspaces of \(C(K)\)}
\author{Run \texttt{fa\_banach\_001}, lane 5}
\date{Archived inside the full packet on 2026-06-19}

\begin{document}
\maketitle

\section*{Historical Status}

This note records the first partial result obtained for the tensor-factor
descent question following Proposition 2.8 of Aron--Choi--Kim--Lee--Martin,
arXiv:1305.6420.  It is superseded by the parent packet, which proves the
qualitative result for every nonzero tensor factor \(E\).

\section*{Theorem}

\begin{theorem}
Let \(K\) be a compact Hausdorff space and let \(E\) be a closed linear
subspace of \(C(K)\) containing \(1_K\).  Let \(X\) and \(Y\) be Banach spaces.
If \((X,E\ot Y)\) has the BPBp with function \(\eta\), then \((X,Y)\) has the
BPBp with the same function \(\eta\).
\end{theorem}

\section*{Proof}

By injectivity of the injective tensor norm, the inclusion \(E\hookrightarrow
C(K)\) induces an isometric embedding
\[
 E\ot Y\hookrightarrow C(K)\ot Y\subset C(K,Y).
\]
Fix \(\varepsilon>0\).  Let \(T\in S_{\mathcal L(X,Y)}\) and \(x_0\in S_X\)
satisfy \(\|Tx_0\|>1-\eta(\varepsilon)\).  Define
\[
 J:Y\to E\ot Y,\qquad J(y)=1_K\otimes y,
\]
and set \(\widetilde T=J\circ T\).  Since \(J\) is an isometry,
\(\widetilde T\) is norm one and
\(\|\widetilde T x_0\|=\|Tx_0\|>1-\eta(\varepsilon)\).

Apply the BPBp for \((X,E\ot Y)\).  There exist
\(\widetilde S\in\mathcal L(X,E\ot Y)\) and \(x_1\in S_X\) such that
\[
 \|\widetilde S\|=\|\widetilde Sx_1\|=1,\qquad
 \|\widetilde S-\widetilde T\|<\varepsilon,\qquad
 \|x_1-x_0\|<\varepsilon .
\]
Viewed in \(C(K,Y)\), the function \(\widetilde Sx_1\) attains its norm at
some \(t_1\in K\).  Let \(Q_{t_1}\) be evaluation at \(t_1\), and set
\[
 S=Q_{t_1}\circ\widetilde S.
\]
Then \(\|S\|\leq 1\) and
\[
 \|Sx_1\|=\|(\widetilde Sx_1)(t_1)\|=1,
\]
so \(S\) attains its norm at \(x_1\).  Finally, \(Q_{t_1}J=\mathrm{Id}_Y\),
hence
\[
 \|S-T\|\leq\|\widetilde S-\widetilde T\|<\varepsilon .
\]
Thus \((X,Y)\) has the BPBp with the same modulus.

\section*{Supersession Note}

The parent full packet replaces the concrete unital \(C(K)\)-subspace
hypothesis by no hypothesis beyond \(E\neq 0\), at the cost of replacing
\(\eta(\varepsilon)\) by \(\eta(\varepsilon/2)\).  The theorem above remains
as a same-modulus corollary.

\end{document}
